\section{Related Work}
\textbf{Failure detection and runtime monitoring.} Failure detection in manipulation is largely treated either as anomaly/out-of-distribution (OOD) detection using statistical models or as supervised/self-supervised learning on high-dimensional data. In the first setting, Eiband \emph{et al.}~\cite{eiband_intuitive_2019} and Willibald \emph{et al.}~\cite{willibald_collaborative_2020} model successful demonstration features and their time-steps as a Gaussian Mixture Model and flag failures via a Mahalanobis-distance threshold to the time-conditioned nominal trajectory; \cite{hierarchical} conditions on end-effector pose rather than time, arguing that interaction dynamics dominate; and \cite{10.5555/3398761.3398779} fits a Hidden Markov Model and detects anomalies online by comparing observation distributions within a sliding window. These methods are effective but share two limitations: an OOD observation does not imply an impending failure, and time-conditioned nominal profiles do not transfer under spatial perturbation, precluding spatial generalization. More recent post-hoc detectors learn distance metrics with time-varying thresholds~\cite{xu_can_2025}, combine a VLM with temporal action consistency~\cite{agia2024unpacking}, or measure action-chunk entropy and observation novelty via random network distillation~\cite{romer2025fiper}. Action-consistency signals are powerful but presuppose action-chunking policies such as diffusion policy~\cite{chi2024diffusionpolicy} or ACT~\cite{Zhao-RSS-23}; SEAL, in contrast, supervises \emph{any} policy, including pre-programmed skills. A parallel line uses VLMs for monitoring and recovery~\cite{sinha_real-time_2024,etukuru_robot_2024,elhafsi_semantic_2023,du_vision-language_2023,zhang_grounding_2023,pacaud_guardian_2025,agia2024unpacking}, but these either fine-tune on failure datasets or query the VLM intermittently due to latency; we use an LLM only at training time and monitor at upwards of 30\,Hz with a lightweight classifier. In the supervised setting, \cite{9636455} trains a deep multimodal classifier on failure examples, Sliwowski and Lee~\cite{sliwowski_conditionnet_2025} predict action phase and flag phase mismatches, and \cite{11245878} selects among detectors by confidence. Their reliance on failure demonstrations makes them unsuitable when such data is unsafe or hard to collect, which is precisely the regime SEAL targets.

\textbf{Mode grounding and precondition learning.} Estimating the semantic mode from observations may be unsupervised (modes discovered) or supervised (modes defined). Unsupervised approaches segment trajectories by geometric reconstruction error~\cite{shi2023waypointbased} or by inverse RL and feature clustering~\cite{hierarchical}. Among supervised approaches, Eiband \emph{et al.}~\cite{eiband_online_2023} combine symbolic and data-driven skill recognition but need hand-specified pre/post-conditions or many per-skill examples; Hegemann \emph{et al.}~\cite{hegemann_learning_2022} ground sensor values with binary predicates fed to a decision tree; and \cite{9811675} composes predicate classifiers with Bayesian state estimation. Purely predicate-based grounding requires an expert to design and train predicates per task; SEAL instead learns the grounding directly and supports both predicate and continuous observations. A related line, initiated by Konidaris \emph{et al.}~\cite{konidaris2018skills}, learns symbolic representations (initiation sets and effect distributions of skills) from interaction data for high-level \emph{planning}; SEAL learns the analogous regions jointly, but from safe demonstrations alone and for a different purpose: real-time localization and failure monitoring during execution rather than offline symbol construction. Finally, \cite{sliwowski_m2r2_2025} learns multimodal representations for temporal action segmentation~\cite{TAS}, and \cite{shah2022value} learns state representations from skill value functions; these are complementary representation learners that could feed a runtime localizer such as ours.

\textbf{Interactive and active imitation learning.} SEAL's querying mechanism belongs to the tradition of robots that solicit targeted help from a human teacher. Chernova and Veloso~\cite{chernova2009confidence} gate autonomous execution on classifier confidence and request demonstrations in low-confidence states; Cakmak and Thomaz~\cite{cakmak2012designing} study which query types make such requests effective for human teachers. The DAgger family formalizes iterative expert correction~\cite{ross2011dagger}, with SafeDAgger~\cite{zhang2017safedagger} and EnsembleDAgger~\cite{menda2019ensembledagger} invoking the expert only where a learned policy is unreliable or an ensemble disagrees. All of these query at the \emph{action} level to improve a control policy, and their queries arise during (possibly unsafe) autonomous rollouts. SEAL transfers the query-when-uncertain principle to the \emph{task} level: queries are short, prescriptive, segmented demonstrations that ground a known automaton, they are generated offline from the calibrated epistemic uncertainty of a GPC rather than from ensemble disagreement, and answering them never requires the robot to act in an uncertain state.

\textbf{Automaton-based sequencing: from funnels to learned grounding.} The formal backbone of our setting is classical: sequential composition~\cite{burridge1999sequential} showed that goal-directed controllers whose attractors lie in the domain of the next controller compose into robust long-horizon behavior (our Assumption~\ref{as:prim} is exactly this condition), and reactive synthesis~\cite{kressgazit2009temporal} showed how temporal-logic specifications compile into correct-by-construction hybrid controllers, assuming the discrete state can be sensed. Temporal Logic Imitation (TLI)~\cite{wang2022temporal} brought this program to learning-from-demonstration: it specifies admissible transitions via an LTL formula and proves that, if each learned dynamical-system policy satisfies mode invariance and goal reachability, any admissible plan is realized robustly; its grounding, however, is a hand-engineered sensor model. GLiDE~\cite{wang2024grounding} closes this gap by \emph{learning} the grounding: it takes valid transitions from an LLM feasibility matrix and learns mode boundaries by perturbing demonstrations and checking, via an outcome oracle and resets, which perturbed trajectories still succeed. Both are important precedents, and we adopt their central abstraction. We differ in \emph{how the grounding is obtained}: SEAL learns it from a few segmented expert demonstrations with minimal post-processing, avoiding (i) perturbation strategies that must be designed per task, (ii) an environment that resets autonomously, and (iii) an outcome oracle. This is a more realistic setting for demonstrators who are proficient in the task but not in model design, and it is the only one of the three that applies when failures are unsafe to induce. We further show (Sec.~\ref{sec:scooping}) that calibrated confidence lets SEAL suppress spurious transitions without incurring switching latency, a property that neither a hand-tuned automaton nor an uncalibrated grounding provides. We do not compare against TAMP frameworks, which plan over tasks and motions rather than reactively sequencing over a fixed nominal plan; when such frameworks incorporate recovery, they typically replan or rely on the anomaly/VLM detectors discussed above.

%%%%%%%%%%%%%%%%%%%%%%%%%%%%%%%%%%%%%%%%%%%%%%%%%%%%%%%%%%%%%%%%%%%%%%%%%%%%%%%%
